% Draft for arXiv (cs.CL / cs.AI). Compile: pdflatex main && bibtex main && pdflatex main && pdflatex main
% DRAFT — numbers are from the prototype runs of 2026-10-13 on branch proto/npc-agent.
% Private GM prep is described, never quoted. Fill the TODOs before submission.
\documentclass[11pt]{article}
\usepackage[margin=1in]{geometry}
\usepackage[T1]{fontenc}
\usepackage[utf8]{inputenc}
\usepackage{lmodern}
\usepackage{microtype}
\usepackage{booktabs}
\usepackage{tabularx}
\usepackage{enumitem}
\usepackage{xcolor}
\usepackage[hidelinks]{hyperref}
\usepackage{natbib}

\newcommand{\todo}[1]{\textcolor{red}{[TODO: #1]}}
\newcommand{\code}[1]{\texttt{#1}}

\title{Playing the NPC from the Record:\\
Knowledge Horizons, Human Rulings, and Backtests\\
for a Language-Model Agent in a Long-Running Tabletop Campaign}

\author{Kostadis Roussos\thanks{\todo{affiliation, contact}}}

\date{Draft, October 2026}

\begin{document}
\maketitle

\begin{abstract}
Can a language model play a non-player character (NPC) in a long-running tabletop
role-playing campaign? The hard part is not combat. It is social play and long-term
planning: the NPC must act on what it knows and wants, and on nothing else.
We describe a series of prototypes, built on 71 sessions of reviewed,
structured session summaries, that give an agent exactly what one NPC could know as
of a given session, which we call a \emph{knowledge horizon}, and then test the agent
against what actually happened at the table.

Our first approach repeated a mistake that a companion survey of grounding documents
for such campaigns had already named. We asked the game master (GM) a question about
every summary \emph{element} that named the NPC, when the question was about
\emph{events}. Rebuilding on the structure the summaries already author, with scenes
as events, bullets as beats and NPC entries as snapshots, cut the questions for one
NPC from 75 to 19 and made them answerable.

Three further lessons came from the GM's rulings:
\begin{itemize}[nosep]
\item Being present at a scene is not knowing everything the GM wrote about it.
\item When the reviewed text's grammar settles a question, the GM should not be asked it.
\item The GM's private preparation, filtered passage by passage into ``hers'' and
``GM only'', is what makes the character play like the GM.
\end{itemize}

In a question-by-question replay of two scenes, the GM judged the agent's answers
the same in substance as the table's in all 11 teacher-forced turns. In a harder
free-running replay, 6 of the 8 turns the GM scored were the same in substance and
the other 2 partly so, with no leaked knowledge. The GM approved all three of the
agent's proposed off-screen moves.

These results cover one NPC and two scenes. We report them as a working method and a
set of design lessons, not as a measurement.
\end{abstract}

\section{Introduction}

A campaign that runs for years accumulates a world. Dozens of named NPCs carry
grudges, secrets and plans across sessions, and the GM must remember what each of them
knows, and when they learned it, every time one of them speaks or acts. That is the
part of the GM's job we set out to support.

The starting suggestion was to ``use the data to create a player'' with a
reinforcement-learning environment such as Gymnasium~\citep{towers2024gymnasium}. We
set that aside, for two reasons:
\begin{itemize}[nosep]
\item No simulator of a tabletop session exists, so there is nothing to roll out
against.
\item In social play, the reward is the GM's judgement.
\end{itemize}
What the data does support is the interface such environments formalise, applied to
the NPC rather than the player. The NPC is an agent with a partial observation (what
it could know), private state (what it wants and hides), and one turn per session.
Its moves are proposed by a model and ratified by the GM.

Three questions organise the work:
\begin{enumerate}[nosep]
\item What could this NPC know, as of this session?
\item Who is the NPC, beyond what the record says?
\item Does an agent given exactly that play the NPC as the GM did, without using
anything it should not know?
\end{enumerate}

The answers are built from reviewed session summaries, never from transcripts. A
separate human-review pipeline turns each recorded session into a structured summary,
which the GM corrects (Section~\ref{sec:data}). The summaries are therefore the
cleanest record of the campaign, and their structure (chapter, scene, bullet, NPC
entry) turns out to matter more than their content.

\paragraph{Contributions.}
\begin{itemize}[nosep]
\item A \emph{knowledge horizon} for an NPC as of a chapter. It is built
deterministically from authored structure and GM rulings: events the NPC physically
attended, minus beats they could not have perceived, with nothing from later
chapters, not even a count.
\item An account of how a survey of grounding documents identified the flaw in our
first design: asking event questions of elements. Also, how the fix followed the
survey's recommendation to use the structure authors already write.
\item A division of labour, ``LLM renders, humans decide''. It is refined by a rule
that emerged from the GM's feedback: when the reviewed text's grammar settles a
question, do not ask the human.
\item A backtest method: freeze the campaign at the chapter before a scene the NPC
attended, then replay the scene question by question. It runs in two modes,
teacher-forced and free-running, and adds a forward turn of off-screen moves that
the GM approves or rejects.
\item First results on one NPC: 11 of 11 teacher-forced turns judged the same in
substance; 8 of 8 scored free-running turns the same or partly; 3 of 3 off-screen
moves approved; no leaks.
\end{itemize}

\section{Setting and data}\label{sec:data}

\paragraph{Campaigns.}
Two campaigns, both run by the author:
\begin{itemize}[nosep]
\item \emph{Out of the Abyss}: 71 sessions, ending in a murder investigation at the
library-fortress of Candlekeep.
\item \emph{Phandelver and Below: The Shattered Obelisk}.
\end{itemize}
Each session is recorded, transcribed, corrected against a hand-maintained record of
transcription errors, and summarised.

\paragraph{Reviewed structured summaries.}
Each summary has the same parts:
\begin{itemize}[nosep]
\item a chapter title;
\item a \code{\#\# Summary} section of prose;
\item a \code{\#\# Scenes} section, each scene with a heading, a one-line \code{\#\#\#\#}
headline (the GM's recap) and bullets;
\item a \code{\#\# Memorable Moments} section of quotations;
\item an \code{\#\# NPCs} section with one entry per NPC who figures in the session.
\end{itemize}
Scenes are numbered \code{NNN.SS} (chapter, scene). The public copy we use has 419
scenes and 6,275 bullets in \emph{Out of the Abyss}; every non-blank scene line falls
in exactly one bullet, its \emph{beat}.

\paragraph{Published dossiers.}
An earlier pipeline links NPC names to summary items and publishes a per-NPC dossier
whose citations resolve into the summaries. We restrict ourselves to the 65 NPCs with
a published dossier, 27 in \emph{Out of the Abyss} and 38 in \emph{Obelisk}. Every
dossier reproduces exactly under our matcher.

\paragraph{The worked NPC.}
Sylvira Savikas is a Great Reader of Candlekeep. She is drawn into the murder of her
lover, the Keeper of Tomes, and hides a demonic son. Her arc is social and secretive,
with no combat to speak of. 19 scenes name her, and she is physically present at two:
\begin{itemize}[nosep]
\item 061.05, an interview in her sickroom with the player character Daz;
\item 067.03, where she helps the party solve a vault's riddles.
\end{itemize}
Between those two scenes the party killed her son (063.05) without her knowledge. The
summary of 067.03 records that she ``does not yet know that the party killed her
son''. That single bullet became the central test case.

\paragraph{A data-quality finding.}
Building the horizon surfaced a summary error. In chapter 63, five attributions to one
NPC belonged to two others; it was filed for correction. \todo{cite issue or omit}.
Checks that read the record closely find errors in it, and the system must defer to
corrections rather than paper over them.

\section{Background: grounding documents for long campaigns}\label{sec:survey}

A companion survey~\citep{roussos2026grounding} studied how to build \emph{grounding
documents} (world state, campaign state, party, planning) for long-running campaigns
from session records. Three of its findings bear directly on this work:
\begin{description}[nosep,leftmargin=1.5em]
\item[Atomisation destroys attribution (F2).] Breaking a record into fragments and
processing each fragment separately loses who did what, and to whom.
\item[Do not undo a rendering step (F10).] Prose that was rendered from structure
should not be parsed back into structure. The structure is the source.
\item[Identity is the precision decision (F8).] Deciding which mention denotes which
entity is the step that needs a human checkpoint, or an authority such as a
registry. It should not be left to similarity.
\end{description}
The survey recommends building on the structure authors already write (its approaches
11 and 12). In the summaries, scenes are events in authored order, bullets are beats,
and NPC entries are snapshots of state.

The repository that hosts this work encodes a related rule: \emph{LLM renders, humans
decide}. Scope, ordering and attribution are precision decisions with a human
checkpoint. Rendering verified structure into prose is what language models do well.
The pattern is ``LLM extracts, human reviews, LLM renders inside that structure'',
never ``LLM extracts, LLM structures, LLM renders''.

\section{The first approach, and the flaw the survey had already named}\label{sec:flaw}

Our first prototype asked the GM a presence question: for each summary item that names
the NPC, was the NPC there? An \emph{item} was any element that matched one of the
NPC's name forms:
\begin{itemize}[nosep]
\item a scene;
\item a paragraph of \code{\#\# Summary} prose;
\item a memorable moment;
\item an \code{\#\# NPCs} entry, hers or another NPC's.
\end{itemize}
For Sylvira that produced 75 questions.

The GM's reaction, reading the answers back, was: ``we are making the same mistake we
made before --- trying to extract events from elements.'' The survey had already named
this mistake, three ways over.

\paragraph{Atomisation (F2).}
One event became several fragments, each asked separately. The interview of 061.05
appears in the summary as one scene, several prose paragraphs, two memorable moments
and an NPC entry. Asking about each separately invited inconsistent answers about a
single occasion, and multiplied the GM's work.

\paragraph{Undoing a rendering step (F10).}
\code{\#\# Summary} prose is a rendering of the scenes. In the private copy of the
summaries, which scene a prose paragraph sits under was decided by a similarity
heuristic during a refactor. Five of the scenes linked to Sylvira were linked only
through such prose. Asking a presence question about a paragraph was asking the GM to
reverse a rendering, and in places to adjudicate a heuristic's placement.

\paragraph{Event questions asked of non-events.}
An NPC entry is a snapshot of state (``absent from the council meeting due to a
lingering illness''), not an occasion anyone could attend. The question had no right
answer there. That showed: the wording of the answers was revised several times
(present / mentioned / unsure / absent; in the scene / not in the scene / can't tell),
because the question itself was malformed for half the items.

\paragraph{The fix.}
The GM cleared every answer, and the prototype was rebuilt on authored structure
(Table~\ref{tab:elements}).
\begin{itemize}[nosep]
\item \emph{Events} are scenes, and only events are asked about.
\item \emph{Beats} are a scene's headline and top-level bullets, carried with the
event.
\item The NPC's own entries are \emph{snapshots}: shown as state, never asked about.
\item Prose, moments and other NPCs' entries are \emph{context}: shown, never asked
about.
\end{itemize}
Matching reuses the dossier linker's rule unchanged:
\begin{itemize}[nosep]
\item exact, case-sensitive and word-bounded;
\item one combined matcher over every published NPC's name forms, longest first;
\item a form owned by two NPCs is withheld;
\item a generic single word is withheld unless the campaign's canon file rules it safe.
\end{itemize}
Sylvira's 75 questions became 19. A companion who travels with the party, whose name
appears in 282 elements, became 91 events inside a 238-scene span that one ruling
covers.

\begin{table}[t]
\centering
\small
\begin{tabularx}{\linewidth}{@{}lXl@{}}
\toprule
Element & What it is & Role in the prototype \\
\midrule
Scene \code{NNN.SS} & an event, in authored order & \textbf{asked}: was the NPC physically present? \\
Scene bullets & its beats & carried with the event into the horizon \\
The NPC's own entry & a snapshot of state & shown as state; never asked \\
Prose, moments, others' entries & a rendering or highlight & shown as context; never asked \\
\bottomrule
\end{tabularx}
\caption{Authored structure, and what is asked of it. Asking event questions only of
events is the survey's recommendation applied to presence.}
\label{tab:elements}
\end{table}

\paragraph{Answers.}
The GM settled three answers:
\begin{itemize}[nosep]
\item \emph{physically present};
\item \emph{mentioned}: named in the scene, but not there;
\item \emph{neither / unsure}.
\end{itemize}
In the GM's words, ``because you will make mistakes it is also necessary for us to
have a unsure or neither present nor mentioned''. Two ruling forms keep the work small.
A \emph{span} covers a companion present throughout a run of scenes. A \emph{default}
answers every remaining event as \emph{mentioned}, for an NPC present at only a few.
Sylvira needed 19 answers; the companion needed one span.

\section{The knowledge horizon}\label{sec:horizon}

The horizon of an NPC at chapter $N$ has two parts.
\begin{description}[nosep]
\item[Part 1: what the NPC could know.] The events answered \emph{physically present}
up to $N$, in order, each with the beats the NPC perceived.
\item[Part 2: for the GM only.] The NPC's latest entry and any events after it, the
beats left out, and the events answered \emph{mentioned}, \emph{neither} or not yet
answered.
\end{description}
Nothing after chapter $N$ appears, not even a count or a first-seen chapter, because a
later chapter number tells an agent the future. An agent is given Part~1 only. The
tool refuses to produce it while any event naming the NPC is unanswered.

\subsection{Being there is not knowing everything the record says}

With Sylvira's answers in, her horizon at chapter 67 correctly placed her son's death
(063.05) outside what she witnessed, and still carried it in. The summaries are
written from the GM's side of the screen. Scene 067.03, which she attends, includes
the beat ``Sylvira is Moziqodo's mother, but does not yet know that the party killed
her son.''

Two consequences were built:
\begin{itemize}[nosep]
\item \emph{Not perceived.} The GM can mark beats of an attended event as not
perceived, at the granularity the summary wrote (a headline or a top-level bullet).
\item \emph{Snapshots are the GM's notes.} The NPC's entries are the GM's record
\emph{about} the NPC. They stay out of Part~1.
\end{itemize}

\subsection{Knowledge notes, and not asking what grammar settles}\label{sec:grammar}

A fixed search found the other beats that say someone does not know something: 21 in
\emph{Out of the Abyss} and 7 in \emph{Obelisk}. It covers phrasings such as ``does not
know'', ``unaware of'' and ``has not been told'', in scene bullets, outside quoted
speech. The same words cover very different things:
\begin{itemize}[nosep]
\item a GM's note on a secret;
\item a remark about what the players do not know;
\item a character telling the party they do not know something.
\end{itemize}

Our review page went through three designs, each corrected by the GM.
\begin{enumerate}[nosep]
\item \textbf{Asking what kind of beat each was.} The GM marked 19 of 21 as a secret
and named nobody, explaining: ``Most of the cards were `the object of the sentence
doesn't know'.'' Every card does say someone doesn't know; the kind was the wrong
question.
\item \textbf{Asking who doesn't know, and whether the beat gives it away.} On a beat
about an NPC who ``didn't know where the secret room was'', the GM asked: ``What am I
saying he doesn't know?'' The card never showed what.
\item \textbf{Showing what.} The GM's next question was decisive: ``I already verified
that all of the sentences are correct, so why ask me?''
\end{enumerate}

The reviewed text's grammar answers the question.
\begin{itemize}[nosep]
\item A \emph{that}-clause, a bare clause, or ``about''/``of'' X states the fact, so the
beat must be hidden from that character. ``does not yet know \emph{that the party
killed her son}'' is one.
\item A question, a name or nothing states nothing hidden, so the beat stays. ``didn't
know \emph{where the secret room was}'' says only that the room's location is
unknown.
\end{itemize}
Who does not know is read from the sentence when the sentence names them. The GM is
asked only when it gives a description instead (``the dragon''). In \emph{Out of the
Abyss}, 5 beats hide something and 16 keep, and the GM answered two questions in
total.

We take this as a refinement of ``humans decide''. The human's review is already in
the text; asking again wastes it and invites inconsistency. A precision decision needs
a human once, and after that it is data.

\section{From horizon to agent}\label{sec:agent}

The agent is given four inputs; Table~\ref{tab:inputs} lists them. Only one module in
the package calls a model, and a static check over the code enforces that. Everything
else is deterministic and refuses rather than guesses.

\begin{table}[t]
\centering
\small
\begin{tabularx}{\linewidth}{@{}lXl@{}}
\toprule
Input & Rule & Decided by \\
\midrule
Persona at $N$ & the NPC's own entries up to $N$, sentence by sentence: keep the descriptive first sentence and sentences she is the subject of; drop reports of her, sentences about others, and any that state what she doesn't know; every drop listed with its reason & code (grammar) \\
Horizon at $N$ & Part 1 of the knowledge horizon & code + GM presence rulings \\
Private sheet & passages of the GM's prep, each ruled \emph{hers} (she knows or believes it) or \emph{GM only}; only hers are shown & GM \\
The scene & its title and one-line headline only; the rest is the answer key & code \\
\bottomrule
\end{tabularx}
\caption{What the agent is given. Scope comes from code and from the GM's rulings;
the model writes only prose.}
\label{tab:inputs}
\end{table}

\paragraph{Prep is written from the GM's side of the screen.}
The GM's session preparation holds what the record lacks. For Sylvira at 061.05 that
was her private reasons and her secret. The record cannot supply these because they
are first written down in the very scene being predicted. But prep also holds things
she must not know: other characters' secrets, mechanics, stage directions, later
scenes. It sometimes holds rulings about her mind that differ from the record.

So code splits a prep file into passages (paragraphs, list items, quotation runs,
table rows). The GM marks each one, and a passage's identifier goes stale when its
text changes, so a ruling is never silently reattached. For the two scenes, 5 of 13
and 2 of 15 passages were ruled hers. The prep is read from a private repository and
never enters the public record or this paper.

\paragraph{Leak check.}
Flags for the GM, never verdicts. A sentence of the agent's answer is flagged when it
carries a hidden beat's content words: all of them, or all but one when there are
three or more, because paraphrase swaps a word. A name is also flagged when the
campaign uses it but the agent was not given it; names first seen after $N$ are marked
as future.

The first real run produced 11 flags, all false positives: ordinary words opening a
quotation (``And'', ``Standing''). The rule now treats a capitalised word the summaries
also use in lower case as an ordinary word. A planted leak (``I know the party killed
Moziqodo'', said to a character she had not met) is still caught. We tuned toward
over-flagging on purpose: a missed leak is the costly error, and the GM can dismiss a
false flag in one click.

\section{Backtests}\label{sec:backtests}

\paragraph{Protocol.}
For each scene the NPC attended:
\begin{itemize}[nosep]
\item freeze the campaign at the chapter before it;
\item give the agent the inputs of Table~\ref{tab:inputs};
\item compare the agent's answers with what happened at the table;
\item have the GM score each answer on a review page.
\end{itemize}
The agent was a frontier large language model \todo{model and version} accessed
through a subscription command-line backend. All prompts and answers are kept with a
hash of their inputs, so runs are reproducible and answers are never silently
regenerated.

\subsection{One turn, no prep: too thin}

Our first backtest asked for one answer of 150 to 300 words per scene. The GM's
verdict was ``it was shorter than the scene\ldots\ Sylvira offers less information.''
There were three causes:
\begin{itemize}[nosep]
\item a length cap not sized to a 33-bullet interview;
\item a single turn in which nobody asked her anything, when at the table every
revelation is an answer to a question;
\item most of what she says at 061.05 is first written down in that scene, so it
existed only in the GM's prep.
\end{itemize}
At 067.03 the agent said the riddles were not her field; the real Sylvira found them
``almost too easy''. Nothing she had been given said she knew Candlekeep's secrets.

\subsection{Question-by-question replay}

A fixed grammatical rule splits the scene into turns. A sentence is hers when its
subject is her: \emph{She}, her name, her possessive, or a subject she shares, after an
optional leading clause such as ``Asked who might want to frame her,''. A turn is a
run of her sentences; what comes between two turns is the next turn's prompt. The
interview of 061.05 becomes 9 turns, and the vault scene 2.

The replay runs in two modes.
\begin{description}[nosep]
\item[Teacher-forced.] The agent sees the real scene so far, including her real
earlier answers.
\item[Free-running.] The agent sees its own earlier answers in place of hers. The
party's lines stay real, since the players cannot be simulated, so they may react to
things this Sylvira never said.
\end{description}

\begin{table}[t]
\centering
\small
\begin{tabular}{@{}llcccc@{}}
\toprule
Scene & Mode & Turns & Same & Partly & Different \\
\midrule
061.05 (interview) & teacher-forced & 9 & 9 & 0 & 0 \\
067.03 (vault)     & teacher-forced & 2 & 2 & 0 & 0 \\
061.05 (interview) & free-running   & 9 (6 scored) & 4 & 2 & 0 \\
067.03 (vault)     & free-running   & 2 & 2 & 0 & 0 \\
\bottomrule
\end{tabular}
\caption{GM judgements of the agent's answers against the table's, turn by turn
(``same substance'' / ``partly'' / ``different''). No leak flags in any answer.}
\label{tab:replay}
\end{table}

\paragraph{Results.}
Table~\ref{tab:replay} gives the scores.
\begin{itemize}[nosep]
\item Teacher-forced, all 11 turns were judged the same in substance.
\item Free-running, 6 of the 8 scored turns were the same and 2 partly so. The GM's
summary: ``This worked really well for Sylvira.''
\item With the clue answers ruled hers, she gives the vault's clues in both modes, as
she did at the table.
\end{itemize}

The divergence in the free run is one of pacing. The GM's prep makes her secret
conditional on being questioned sincerely. The agent applied that condition at the
first moment it was met, the turn after the summary records that Daz ``persuades her of his sincerity''. At
the table the same secret came out four turns later. The two partly scores fall
exactly there. Once she has confessed, her later answers match the table again,
reacting to her own earlier confession.

This exposes something prep does not say: a secret's release condition can be stricter
in play than on paper.

\subsection{Forward turn}

Frozen at chapter 66, before her next attended scene, the agent was asked what
Sylvira does off-screen before she next meets the party: one to three moves, each with
a reason. Its inputs were her persona, horizon and private sheet, plus the chapter's
summary prose when she attended an event in it.

The GM's plan for her next session was shown to the GM beside the moves, for
comparison only, and never to the agent. All three moves were approved:
\begin{itemize}[nosep]
\item a return to her son;
\item burning her treatments faster than is wise to keep her mind clear;
\item setting her familiar to watch for her impersonator.
\end{itemize}
The first is instructive. Her son was already dead, killed by the party at 063.05
without her knowledge. A horizon-bounded agent proposes an errand to him, which is
correct for what she knows and dramatic for the GM. The approvals are stored with the
move text and the prompt's hash, as the first entries of a preference record.

\section{Discussion}\label{sec:discussion}

\paragraph{Ask questions of the unit the question is about.}
The flaw of Section~\ref{sec:flaw} was not a bug in matching. It was a category error,
an event question asked of elements, and the survey's findings predicted its symptoms
exactly: fragmented answers, reversed renderings, and wording that churned because
there was no right answer. The fix was to stop deriving structure and use the
structure the summaries already author.

\paragraph{Do not ask what the text already settles.}
Each correction the GM made to our review pages moved a decision from the human to
the reviewed text: the kind of a beat, what it says is unknown, who does not know it.
The human checkpoint belongs where the text is silent or ambiguous: a description
instead of a name, presence at a scene, which prep passages are hers. It does not
belong at every decision.

\paragraph{The record is not enough; the prep is not safe.}
The record gives an agent what an NPC experienced. It does not give who they are or
what they hide, because that is first written in the scene a backtest is trying to
predict. The GM's prep gives it, but mixed with everything the NPC must not know. A
per-passage human ruling was the cheapest safe filter we found: 28 decisions across
two scenes.

\paragraph{Limitations.}
\begin{itemize}[nosep]
\item One NPC, two scenes, 11 turns and 3 moves. This is a demonstration of a method,
not a measurement.
\item The scorer is the GM, who is also the author and whose prep the agent read.
Blind and multi-rater scoring are needed.
\item Turn splitting is grammatical and makes visible mistakes. The free-running
replay keeps the party's lines fixed, so drift partly reflects the test format, not
the agent.
\item The leak check compares words and names, not meanings, so a paraphrase that
shares too few content words with a hidden beat can pass.
\item The forward turn was judged, not checked: this NPC has no written rules to test
moves against.
\end{itemize}

\paragraph{Future work.}
\begin{itemize}[nosep]
\item An NPC whose GM has written guardrails, about 50 checks for one antagonist,
to test forward moves mechanically before the GM sees them.
\item More NPCs, to learn whether this generalises.
\item A way to give a prep secret a release condition.
\item With enough approve/reject pairs, preference learning over moves. Here the
Gymnasium-style interface becomes useful again, with the GM's ratifications as the
signal rather than a simulated reward.
\end{itemize}

\section{Related work}
\todo{Expand.} Generative agents simulate believable behaviour from memory streams and
reflection~\citep{park2023generative}; our setting differs in having a canonical
record and a human authority over what a character knows. Role-play with language
models~\citep{shanahan2023roleplay} frames the model as playing characters rather than
being one; our horizon is an explicit, auditable bound on the character's knowledge.
Gymnasium~\citep{towers2024gymnasium} and PettingZoo~\citep{terry2021pettingzoo}
standardise agent--environment interfaces; we use the interface's vocabulary without a
simulator. The grounding-document survey~\citep{roussos2026grounding} supplies the
design principles this work tests.

\section{Conclusion}
An NPC agent for a long-running campaign can be grounded in reviewed session summaries
with little human effort. The effort goes where the reviewed text is silent:
\begin{itemize}[nosep]
\item 19 presence answers;
\item 2 identity rulings;
\item 28 prep-passage rulings.
\end{itemize}
Everything else is read from authored structure and grammar.

Our first design failed by asking event questions of elements. A survey of grounding
documents had already named that mistake, and the fix was its own recommendation. The
resulting agent played one secretive NPC with the same substance as the GM in every
teacher-forced turn and most free-running ones, without using what she could not
know.

\bibliographystyle{plainnat}
\begin{thebibliography}{9}

\bibitem[Park et~al.(2023)]{park2023generative}
J.~S. Park, J.~C. O'Brien, C.~J. Cai, M.~R. Morris, P.~Liang, and M.~S. Bernstein.
\newblock Generative agents: Interactive simulacra of human behavior.
\newblock In \emph{Proceedings of UIST}, 2023.

\bibitem[Roussos(2026)]{roussos2026grounding}
K.~Roussos.
\newblock Grounding documents for long-running tabletop campaigns.
\newblock Manuscript, October 2026. \todo{arXiv id or URL}.

\bibitem[Shanahan et~al.(2023)]{shanahan2023roleplay}
M.~Shanahan, K.~McDonell, and L.~Reynolds.
\newblock Role play with large language models.
\newblock \emph{Nature}, 623:493--498, 2023.

\bibitem[Terry et~al.(2021)]{terry2021pettingzoo}
J.~K. Terry et~al.
\newblock PettingZoo: Gym for multi-agent reinforcement learning.
\newblock In \emph{Advances in Neural Information Processing Systems}, 2021.

\bibitem[Towers et~al.(2024)]{towers2024gymnasium}
M.~Towers et~al.
\newblock Gymnasium: A standard interface for reinforcement learning environments.
\newblock arXiv:2407.17032, 2024.

\end{thebibliography}

\end{document}
